% =====================================================================
%  INF-5103 -- Concepts Statistiques pour la Science des Donnees
%  Devoir I -- Analyse statistique de quatre jeux de donnees
%
%  Fichier autonome : la classe uqo-memoire n'est PAS utilisee (elle est
%  faite pour un memoire de 100 pages et en consommait 10 a vide). Seul
%  son STYLE est repris : page de garde officielle, marges, en-tete
%  courant filete, conventions typographiques francaises.
%
%  CONTRAINTE : SIX PAGES MAXIMUM, page de garde comprise.
%
%  Compilation :  tectonic rapport.tex     (ou pdflatex deux fois)
%  Figures en PDF ou PNG -- pas d'EPS, tectonic ne les lit pas.
% =====================================================================
% [IA] Code généré par Claude : police 11pt (au lieu de 12pt) pour respecter la limite de 6 pages
\documentclass[11pt,letterpaper]{article}

\usepackage[utf8]{inputenc}
\usepackage[T1]{fontenc}
\usepackage[french]{babel}
\usepackage{lmodern}
% [IA] Code généré par Claude : marges de 2 cm (au lieu de 3 et 2,5 cm) pour respecter la limite de 6 pages
\usepackage[margin=2cm]{geometry}
\usepackage{graphicx}
\usepackage{booktabs}
\usepackage{amsmath}
\usepackage{xcolor}
\usepackage{float}
\usepackage{caption}
\usepackage{fancyhdr}
\usepackage{wrapfig} % [IA] Code généré par Claude : texte à côté d'une figure (inutilisé)
\usepackage[hidelinks]{hyperref}

% --- Style repris du modele UQO --------------------------------------

% En-tete courant filete : titre de section a gauche, page a droite.
\pagestyle{fancy}
\fancyhf{}
\fancyhead[L]{\small\nouppercase{\rightmark}}
\fancyhead[R]{\small\thepage}
\renewcommand{\headrulewidth}{0.6pt}
\renewcommand{\footrulewidth}{0pt}
\renewcommand{\sectionmark}[1]{\markright{\thesection~--~#1}}

% Pas de veuves ni d'orphelines (comme la classe du memoire).
\widowpenalty=10000
\clubpenalty=10000

\captionsetup{font=small,labelfont=bf,skip=4pt}

% [IA] Code généré par Claude : espacements resserrés (titres, flottants, paragraphes) pour respecter la limite de 6 pages
\usepackage{titlesec}
\titlespacing*{\section}{0pt}{1.6ex plus .4ex minus .2ex}{.8ex}
\titlespacing*{\subsection}{0pt}{1.2ex plus .3ex minus .2ex}{.5ex}
\titleformat{\subsubsection}[runin]{\normalfont\normalsize\bfseries}{\thesubsubsection}{.5em}{}[.]
\titlespacing*{\subsubsection}{0pt}{.8ex plus .2ex minus .2ex}{.5em}
\titlespacing*{\paragraph}{0pt}{1ex plus .2ex minus .2ex}{1em}
\setlength{\intextsep}{6pt plus 2pt minus 2pt}
\setlength{\textfloatsep}{8pt plus 2pt minus 2pt}
\setlength{\floatsep}{6pt plus 2pt minus 2pt}
\setlength{\parskip}{0pt}
\captionsetup{skip=3pt}

% --- Outils de redaction ---------------------------------------------

% Texte temporaire, bien visible : AUCUN placeholder ne doit survivre a
% la remise.   Verification :  grep -c 'ph{' rapport.tex   ->  0
\newcommand{\ph}[1]{{\color{gray}\itshape [#1]}}

% Intitule de question en retrait de paragraphe : compact, et chaque
% question de l'enonce reste reperable d'un coup d'oeil.
\newcommand{\qst}[1]{\paragraph{#1.}}

% Inclut la figure si elle existe, sinon dessine un cadre temoin : le
% document compile avant meme que les figures soient produites.
\newcommand{\figplace}[2][0.62\textwidth]{%
  \IfFileExists{#2}{\includegraphics[width=#1]{#2}}%
  {\fbox{\parbox[c][0.32\textwidth][c]{#1}{\centering\ttfamily\small
   \detokenize{#2}\\[2pt]\normalfont\rmfamily (figure a produire)}}}%
}

% --- Chiffres et tableaux produits par le notebook -------------------
% Le notebook ecrit resultats/valeurs.tex ; chaque \val{cle} y pioche
% son chiffre. Aucun nombre n'est recopie a la main dans ce fichier :
% il suffit de reexecuter le notebook pour que le rapport suive.
% Une cle absente s'affiche en rouge dans le PDF au lieu de casser la
% compilation, donc on repere tout de suite ce qui manque.
\newcommand{\val}[1]{%
  \ifcsname val@#1\endcsname
    \csname val@#1\endcsname
  \else
    {\color{red}\textbf{[?#1]}}%
  \fi}
\IfFileExists{resultats/valeurs.tex}{% Généré par le notebook -- ne pas éditer à la main.
\expandafter\gdef\csname val@jeu1-ecart\endcsname{25,3\,points}
\expandafter\gdef\csname val@jeu1-n-a\endcsname{1200}
\expandafter\gdef\csname val@jeu1-n-b\endcsname{1200}
\expandafter\gdef\csname val@jeu1-taux-a\endcsname{70,3\,\%}
\expandafter\gdef\csname val@jeu1-taux-b\endcsname{45,1\,\%}
}{}

% Tableau genere par le notebook (pandas .to_latex), meme principe :
% \tbl{nom} charge resultats/nom.tex, quel que soit le nom.
\newcommand{\tbl}[1]{%
  \IfFileExists{resultats/#1.tex}{\input{resultats/#1.tex}}%
  {\fbox{\parbox{0.8\textwidth}{\centering\ttfamily\small
   \detokenize{resultats/#1.tex}\\[2pt]\normalfont\rmfamily
   (tableau a produire par le notebook)}}}}

% --- A personnaliser -------------------------------------------------
\newcommand{\etudiant}{Alex D'Andrea}
\newcommand{\matricule}{DANA21079300}
\newcommand{\titredev}{Analyse statistique de quatre jeux de données}
\newcommand{\coursdev}{INF-5103 -- Concepts statistiques pour la science des données}
\newcommand{\enseignant}{Mohamed Lamine Allaoui}
\newcommand{\datedev}{9 octobre 2026}
% ---------------------------------------------------------------------

\begin{document}

% =====================================================================
%  PAGE DE GARDE  (mise en page du modele de memoire de l'UQO)
% =====================================================================
\begin{titlepage}
\thispagestyle{empty}
\centering
\vspace*{\fill}

{\MakeUppercase{\bfseries Université du Québec en Outaouais}\par}
\vspace{0.3cm}
{Département d'informatique et d'ingénierie\par}

\vspace{1.5cm}
{\large\scshape Devoir I\par}

\vspace{0.8cm}
\rule{\linewidth}{0.8pt}\\[0.5cm]
{\large\MakeUppercase{\titredev}\par}
\vspace{0.5cm}
\rule{\linewidth}{0.8pt}

\vspace{1.5cm}
TRAVAIL PRÉSENTÉ\par
DANS LE CADRE DU COURS\par
\MakeUppercase{\coursdev}

\vspace{2cm}
PAR\par
\vspace{0.4cm}
{\MakeUppercase{\etudiant}\par}
{Matricule : \matricule\par}

\vspace{1.5cm}
{Travail individuel remis à\par \enseignant\par}

\vspace{1.5cm}
\MakeUppercase{\datedev}

\vspace*{\fill}
\end{titlepage}

% =====================================================================
\section{Plateforme d'apprentissage en ligne}
% Q1-Q5 : taux global, taux par niveau, paradoxe de Simpson,
%         standardisation par les effectifs du groupe entier, autre variable.

\subsection{Taux de réussite global (Q1)}


% [IA] Code généré par Claude : mise en page texte + figure côte à côte (minipage, \figplace, \captionof)
\noindent
\begin{minipage}[c]{0.55\textwidth}

% Cles disponibles pour cette question, produites par le notebook :
%   \val{jeu1-n-a}  \val{jeu1-n-b}  \val{jeu1-taux-a}  \val{jeu1-taux-b}
%   \val{jeu1-ecart}
% Les inserer dans la redaction plutot que de retaper les chiffres.
La figure suivante présente le taux de réussite global pour les deux méthodes.
Il est donc possible que la méthode A ait un meilleur taux de réussite global.
Si on regarde uniquement cette information, la méthode A serait à recommander.
\end{minipage}\hfill
\begin{minipage}[c]{0.42\textwidth}
  \centering
  \figplace[0.8\linewidth]{figures/jeu1_reussite.pdf}
  \captionof{figure}{Taux de réussite global.}
  \label{fig:jeu1-reussite}
\end{minipage}

\subsection{Taux par niveau initial (Q2)}

% [IA] Code généré par Claude : tableau et figure côte à côte (Q2)
\noindent
\begin{minipage}[c]{0.44\textwidth}
  \centering\small
  \captionof{table}{Résultats par niveau initial.}
  \label{tab:jeu1-niveau}
  \tbl{jeu1-Q1}
\end{minipage}\hfill
\begin{minipage}[c]{0.54\textwidth}
  \centering
  \figplace[0.8\linewidth]{figures/jeu1_niveau.pdf}
  \captionof{figure}{Résultats par niveau initial.}
  \label{fig:jeu1-niveau}
\end{minipage}
En examinant le niveau initial, on constate que la méthode B est testée sur un bien plus grand nombre d'étudiants débutants que la méthode A.
On constate que la méthode B est plus de deux fois plus efficace que la méthode A pour les étudiants débutants.
En fait, elle est plus efficace pour tous les niveaux.
Avec cette analyse, on recommande donc la méthode B.

\subsection{Confrontation des questions 1 et 2 (Q3)}
Les deux méthodes n'ont pas le même nombre d'étudiants par niveau initial.
On constate que la méthode A est testée sur un bien plus grand nombre d'étudiants avancés que la méthode B.
Cela donne un plus grand nombre de réussites pour la méthode A, mais ce n'est pas représentatif de l'efficacité de la méthode.
Vu que la méthode B est plus efficace pour tous les niveaux, on recommande donc la méthode B.

\subsection{Taux standardisés (Q4)}
\begin{minipage}[c]{0.50\textwidth}

% Cles disponibles pour cette question, produites par le notebook :
%   \val{jeu1-n-a}  \val{jeu1-n-b}  \val{jeu1-taux-a}  \val{jeu1-taux-b}
%   \val{jeu1-ecart}
% Les inserer dans la redaction plutot que de retaper les chiffres.
La figure suivante présente le taux de réussite global pondéré pour les deux méthodes.
On constate que la méthode B est plus efficace que la méthode A.
On constate une quasi-inversion de la figure 1.
Avec la pondération, on recommande donc la méthode B.
\end{minipage}\hfill
\begin{minipage}[c]{0.45\textwidth}
  \centering
  \figplace[0.8\linewidth]{figures/jeu1_reussite_ponderee.pdf}
  \captionof{figure}{Taux de réussite global pondéré.}
  \label{fig:jeu1-reussite-pondéré}
\end{minipage}

\subsection{Reprise avec une autre variable (Q5)}
\begin{minipage}[c]{0.45\textwidth}

La figure suivante présente les taux de réussite par région.
Chaque région a un nombre similaire d'étudiants pour les deux méthodes.
On constate que pour toutes les régions, la méthode A a un meilleur taux de réussite que la méthode B.
Il est probable que ceci soit dû aux régions ayant une répartition similaire pour les niveaux initiaux, et qui n'ont donc pas d'impact sur le taux de réussite.
Une comparaison avec les niveaux initiaux par région serait nécessaire pour confirmer cette hypothèse.
\end{minipage}\hfill
\begin{minipage}[c]{0.55\textwidth}
  \centering
  \figplace[0.8\linewidth]{figures/jeu1_region.pdf}
  \captionof{figure}{Taux de réussite par région.}
  \label{fig:jeu1-region}
\end{minipage}
% =====================================================================
\section{Consommation électrique d'un immeuble de bureaux}
% Q1-Q6 : correlation, nuage de points, correlation par groupe de
%         temperature, minimum, occupation/jour_ferie, lecon sur r.

\subsection{Corrélation température--consommation (Q1)}

Le coefficient de corrélation de la température et de la consommation est de $r = \val{correlation-coef}$.
Ce coefficient est positif, mais très faible, indiquant une corrélation positive quasi nulle.

\subsection{Nuage de points (Q2)}
\begin{minipage}[c]{0.45\textwidth}

Le nuage de points de la figure \ref{fig:jeu2-correlation} montre la relation entre la température et la consommation.
On constate que la consommation d'énergie diminue avec l'augmentation de la température, jusqu'à environ 17 degrés Celsius.
À partir de ce point, la consommation d'énergie augmente avec l'augmentation de la température.
Ceci semble refléter le fait que l'immeuble est chauffé lorsque la température est basse, et refroidi lorsque la température est élevée.
\end{minipage}\hfill
\begin{minipage}[c]{0.55\textwidth}
  \centering
  \figplace[0.8\linewidth]{figures/jeu2_nuage_conso_temperature.pdf}
  \captionof{figure}{Nuage de points : consommation en fonction de la température.}
  \label{fig:jeu2-correlation}
\end{minipage}


\subsection{Corrélation de part et d'autre d'un seuil (Q3)}
\begin{minipage}[c]{0.55\textwidth}

Pour trouver le seuil de température qui maximise la corrélation, nous avons calculé les coefficients de corrélation pour des seuils variant de 15 à 20 degrés Celsius.
Nous avons fait le total des coefficients de corrélation pour les températures inférieures et supérieures à chaque seuil, en inversant le signe du coefficient pour les températures inférieures.
Le seuil de température qui maximise la corrélation est de \val{seuil-optimal} degrés Celsius.
Le coefficient de corrélation pour les températures inférieures à ce seuil est de \val{coef-froid}, et celui pour les températures supérieures à ce seuil est de \val{coef-chaud}.
La figure \ref{fig:jeu2_coefs_par_seuil} montre ces coefficients pour différents seuils.
\end{minipage}\hfill
\begin{minipage}[c]{0.40\textwidth}
  \centering
  \figplace[0.8\linewidth]{figures/jeu2_coefs_par_seuil.pdf}
  \captionof{figure}{Coefficients de corrélation par seuil de température.}
  \label{fig:jeu2_coefs_par_seuiln}
\end{minipage}

\subsection{Température de consommation minimale (Q4)}

Nous avons ensuite trouvé la température avec la consommation minimale en calculant la moyenne de la consommation pour chaque température.
La température avec la consommation minimale est de \val{jeu2-temp-min} degrés Celsius.

\subsection{Rôle de l'occupation et des jours fériés (Q5)}
% [IA] Code généré par Claude : légende, étiquette et mise en forme du tableau des corrélations (Q5)
\begin{table}[H]
  \centering\small
  \caption{Coefficients de corrélation entre consommation, occupation,
    journée de semaine non fériée et température.}
  \label{tab:jeu2-correlations}
  \tbl{jeu2-Q5-correlations}
\end{table}
Le tableau \ref{tab:jeu2-correlations} montre les coefficients de corrélation entre la consommation, l'occupation, la journée de semaine non fériée et la température.
On constate que les jours de semaine non fériés et l'occupation ont une corrélation proche de 1.
Ils ont également une corrélation très similaire avec la consommation.
On peut donc fortement soupçonner qu'ils sont colinéaires et affectent la consommation de la même manière.
On constate également que les deux facteurs ont une corrélation quasi nulle avec la température, indiquant un risque faible de colinéarité avec la température.
Ceci est important pour s'assurer que les facteurs n'affectent pas la consommation de la même manière que la température, et donc ne faussent pas l'analyse de la corrélation entre la consommation et la température.

 

\subsection{Ce que ce jeu apprend sur le coefficient de corrélation (Q6)}
Cette section démontre que plusieurs choses peuvent affecter une variable.
Il est donc important de valider ce qui affecte quoi, et comment, pour s'assurer que les corrélations sont bien interprétées.
Deux variables qui affectent une troisième variable peuvent avoir des corrélations très différentes avec cette variable, mais ne sont pas nécessairement corrélées entre elles.

% =====================================================================
\section{Dossiers de patients d'une clinique}
% Q1-Q6 : description et anomalies, nettoyage documente, moyenne/mediane
%         de l'age, correlation age-tension avant/apres, IMC, tabagisme.

\subsection{Description et anomalies détectées (Q1)}

% [IA] Code généré par Claude : insertion du tableau de description des dossiers de patients
\begin{table}[H]
  \centering\small
  \caption{Description statistique des colonnes numériques des dossiers de
    patients (données brutes).}
  \label{tab:jeu3-description}
  \tbl{jeu3-Q1-description}
\end{table}
Le tableau \ref{tab:jeu3-description} montre la description statistique des colonnes numériques des dossiers de patients.
On constate que l'écart-type de l'âge est impossible, ainsi que son maximum.
Fort probablement, il y a des valeurs de 999 pour l'âge, ce qui est une valeur aberrante.
On constate également que la tension a un minimum de $-1$, ce qui est impossible.
Ensuite, on constate que la taille minimale est de 1,5.
Ceci est impossible et doit être investigué.
Le traitement de ces anomalies est décrit dans la section suivante.
Le tabagisme et le sexe sont des variables catégorielles et ne sont pas décrits dans le tableau.
Ils comportent uniquement les mêmes valeurs que celles décrites dans l'énoncé et ne comportent pas d'anomalies.


\subsection{Préparation des données (Q2)}
Pour toutes les anomalies, les détails des valeurs peuvent être trouvés dans le notebook.
La technique de l'IQR a été utilisée pour détecter les valeurs aberrantes.
Après nettoyage, 771 des 800 patients demeurent dans l'analyse.
\subsubsection{Âge}
L'âge et la tension ont eu le même processus de nettoyage.
Pour l'âge, nous avons regardé les voisins les plus proches, avec moins de 5 kg de différence (de même sexe et de même statut tabagique).
Avec ces voisins, nous avons calculé la distance standardisée avec les voisins.
Nous avons ensuite fait une moyenne pondérée des âges des voisins.
Les patients avec des voisins trop éloignés ont été exclus. 
Au total, 15 patients ont été exclus pour l'âge.
25 patients ont eu leur âge imputé par la moyenne pondérée de leurs voisins.
\subsubsection{Tension}
Pour la tension, nous avons regardé les voisins les plus proches, avec moins de 5 ans de différence (de même sexe et de même statut tabagique).
Avec ces voisins, nous avons calculé la distance standardisée avec les voisins, comme pour l'âge.
Nous avons ensuite fait une moyenne pondérée des tensions des voisins.
Les patients avec des voisins trop éloignés ont été exclus.
Au total, 14 patients ont été exclus pour la tension.
11 patients ont eu leur tension imputée par la moyenne pondérée de leurs voisins.
\subsubsection{Taille}
La taille comportait des valeurs supérieures et inférieures aux bornes de l'IQR.
Les valeurs supérieures étaient à moins d'un cm de la borne. 
Elles ont donc été incluses dans l'analyse.
Les valeurs inférieures étaient entre 1,54 et 1,92. 
Si on suppose que les valeurs étaient en m au lieu de cm, alors les valeurs sont entre 154 et 192 cm, ce qui est plausible.
Nous avons donc multiplié les valeurs inférieures par 100 pour les mettre en cm.
\subsubsection{Poids}
Les valeurs aberrantes en poids étaient toutes supérieures à la borne de l'IQR.
Les valeurs sont à moins de 5 kg de la borne et sont donc plausibles.
Elles ont donc été incluses dans l'analyse.

\subsection{Âge : moyenne et médiane, avant et après (Q3)}
Avant le nettoyage, l'âge moyen était de \val{jeu3-age-moyenne-brut} ans, et la médiane était de \val{jeu3-age-mediane-brut} ans.
Après le nettoyage, l'âge moyen est de \val{jeu3-age-moyenne-corrige} ans, et la médiane est de \val{jeu3-age-mediane-corrige} ans.
On constate que la moyenne a grandement diminué, mais la médiane a peu changé.
En fait, la moyenne est plus proche de la médiane après le nettoyage, indiquant une meilleure distribution de l'âge après le nettoyage.

\subsection{Corrélation âge--tension systolique, avant et après (Q4)}
Regardons la corrélation avant le nettoyage.
Le coefficient de corrélation de l'âge et de la tension systolique est de $r = \val{jeu3-corr-age-tension-brut}$.
Ce nombre est très faible, indiquant une corrélation quasi nulle.
Cela est probablement dû aux valeurs aberrantes dans l'âge et la tension systolique.
Après le nettoyage, le coefficient de corrélation de l'âge et de la tension systolique est de $r = \val{jeu3-corr-age-tension-corrige}$.
Cette valeur est beaucoup plus élevée, indiquant une corrélation significative entre l'âge et la tension systolique après le nettoyage.
Cela reflète des faits connus, et est donc plus réaliste.

\subsection{Indice de masse corporelle (Q5)}
% [IA] Code généré par Claude : insertion du tableau de la tension moyenne par classe d'IMC
\begin{table}[H]
  \centering\small
  \caption{Tension systolique moyenne par classe d'IMC (classes de Santé Canada).}
  \label{tab:jeu3-imc}
  \tbl{jeu3-Q5-imc}
\end{table}
On peut voir que la tension systolique moyenne augmente avec l'IMC.
La classe III est omise, car elle ne comporte aucun patient.


\subsection{Tabagisme et sexe (Q6)}

% [IA] Code généré par Claude : insertion du tableau des corrélations de la tension avec le sexe et le tabagisme
\begin{table}[H]
  \centering\small
  \caption{Corrélation entre la tension systolique et des indicateurs 0/1
    du sexe et de chaque statut tabagique (données corrigées).}
  \label{tab:jeu3-correlations-tension}
  \tbl{jeu3-Q6-correlations-tension}
\end{table}
Le tableau ne présente aucune corrélation entre la tension et le sexe.
Le tableau présente une faible corrélation inverse entre le fait d'être non-fumeur et la tension.
Finalement, le tableau présente une faible corrélation positive entre le fait d'être fumeur et la tension.

% =====================================================================
\section{Temps de réponse d'un service web}
% Q1-Q4 : moyenne/mediane/p95, histogrammes lin et log, explication de
%         la forme, choix d'un indicateur unique.

\subsection{Moyenne, médiane et 95\textsuperscript{e} centile (Q1)}
% [IA] Code généré par Claude : insertion du tableau des statistiques de latence (jeu 4, Q1)
\begin{table}[H]
  \centering\small
  \caption{Moyenne, médiane et 95\textsuperscript{e} centile de la latence,
    et nombre de requêtes à moins de 10\,\% de la moyenne.}
  \label{tab:jeu4-latence}
  \tbl{jeu4-Q1-latence}
\end{table}
Avec la moyenne et la médiane si différentes, on constate que la distribution est asymétrique.
Comme aucune valeur n'est près de la moyenne, on suppose que toutes les données seront centrées sur la médiane.
Drôlement, le 95\textsuperscript{e} centile est loin de la médiane, indiquant une longue queue de valeurs élevées, mais peu de ces valeurs.


\subsection{Histogrammes en échelle linéaire et logarithmique (Q2)}

\begin{figure}[H]
  \centering
  \figplace[0.36\textwidth]{figures/jeu4_hist_lin.pdf}\hfill
  \figplace[0.36\textwidth]{figures/jeu4_hist_log.pdf}
  \caption{Distribution des latences : échelle linéaire (gauche) et
  logarithmique (droite).}
  \label{fig:jeu4-hist}
\end{figure}
Les figures représentent chacune une façon différente de visualiser la distribution des latences.
Premièrement, la figure linéaire, à gauche, montre effectivement une queue à droite, avec un énorme pic à gauche.
Cela montre une distribution asymétrique à droite, mais il est difficile de savoir combien de données se trouvent à droite.
Ensuite, sur la figure logarithmique, on constate deux pics, un à gauche et un plus petit à droite.


\subsection{Origine de la forme observée (Q3)}
% [IA] Code généré par Claude : insertion du tableau de la latence par version du service (jeu 4, Q3)
\begin{table}[H]
  \centering\small
  \caption{Latence par version du service : effectif, moyenne, médiane et
    écart-type.}
  \label{tab:jeu4-version}
  \tbl{jeu4-Q3-version}
\end{table}



% [IA] Code généré par Claude : texte et figure des histogrammes par version côte à côte (minipage, \figplace, \captionof)
\noindent
\begin{minipage}[c]{0.38\textwidth}

  La table \ref{tab:jeu4-version} montre la latence par version du service.
On constate que la version 2.4.0 a une latence moyenne et médiane beaucoup plus élevée que 2.3.1.
On constate également que la version 2.4.0 a un écart-type beaucoup plus élevé que 2.3.1.
Cela peut expliquer la forme.  

  La figure \ref{fig:jeu4-version} montre la distribution des latences par version du service.
On constate que les deux versions ont deux pics, mais le pic plus lent pour la version 2.4.0 est beaucoup plus grand que celui de la version 2.3.1.
\end{minipage}\hfill
\begin{minipage}[c]{0.60\textwidth}
  \centering
  \figplace[\linewidth]{figures/jeu4_hist_version.pdf}
  \captionof{figure}{Distribution des latences par version du service : échelle
    linéaire (haut) et logarithmique (bas). Axes des ordonnées propres à
    chaque panneau.}
  \label{fig:jeu4-version}
\end{minipage}


\subsection{Indicateur unique à mettre sous surveillance (Q4)}
Après avoir regardé plusieurs indicateurs, la version du service est le meilleur indicateur à surveiller.
La nouvelle version a un impact négatif sur la latence et doit être surveillée pour éviter de dégrader l'expérience utilisateur.
Lorsque plus de données seront disponibles, il sera possible de confirmer ceci.

% =====================================================================
%  Categorie 2 de l'enonce : declarer l'outil et ce qu'il a produit.
%  Une declaration honnete n'entraine aucune penalite.
% =====================================================================
\section*{Déclaration d'utilisation de l'intelligence artificielle}
\markright{Déclaration d'utilisation de l'intelligence artificielle}

L'IA générative Claude a été utilisée pour ce devoir.
L'outil a principalement été utilisé pour la production de gabarits pour le document LaTeX du rapport, ainsi que pour la relecture.
L'outil a préparé les fonctions qui servent à exporter les résultats pour une importation rapide dans le rapport .tex.

% [IA] Lien GitHub ajouté par Claude à la demande de l'étudiant
L'utilisation détaillée de l'IA peut être consultée dans les commits du dépôt git
(\url{https://github.com/Dandalex/Devoir1ConceptStat}).
Chaque étape produite par l'outil est dans son propre commit, avec la signature à cet effet.

\end{document}
